\section{账本事件的区域制度深描：three-kingdoms}

\label{sec:ledger-depth-three-kingdoms}

本组叙事以本卷事件账本的可核年序为骨架，将政治节点放回地方资源、制度中介、人口流动与材料类型中。每条来源能够证明的范围不同，故不以朝廷命令、战报数字或后出传记替代基层社会的完整经验。

\subsection{地方资源与制度接口}

账本条目\texttt{TK-0220-WEI-LUOYANG}记载，黄初元年（220年），曹魏发生“曹丕定都洛阳，以魏的尚书、禁军和北方仓储网络承接汉末朝廷”。这一节点可放在受禅后的新政权需要脱离许县的战时行营格局，并把控制华北的军政机构置于旧东汉都城的条件下理解；其后续影响被账本概括为洛阳成为魏及其后西晋的政治枢纽；都城选择将中原漕运、宫廷警卫和北方士族任用重新连结。就地方尺度而言，事件并不只属于朝廷或统帅：征发、道路、仓储、户籍、市场、寺院与家庭关系都可能决定命令如何抵达、战事如何延长，或接收何以在不同地区呈现不一样的速度。

本条列举的核心材料为《三国志·文帝纪》；《三国志·明帝纪》；《资治通鉴》魏纪。这些材料较适合钉住事件的发生、官方表述或特定人物、地点，却不自动构成对全境人口、收入、兵数、信仰或动机的调查。账本的置信与材料提示是“定都事实可靠；汉魏洛阳城的宫城复用、营建节奏与人口规模须由遗址分期核验”；来源限度进一步说明“魏初都城、汉魏洛阳城与北方行政网络研究”。因此，不能把条目中的政权名称、行政分类或战报修辞直接转译为固定族群和统一社会意志，也不能将制度宣布写成各地同步实施。

在叙事上，更稳妥的做法是把这一事件视为一条需要地方中介才能运作的链：中央的命令、地方官与精英的选择、运输与劳力的条件、以及普通家庭可能面对的迁徙、赋役或安全风险彼此相连。现有材料能够显示其中若干环节，却保留了大量沉默。承认这种沉默并非放弃解释，而是避免以一条编年记录覆盖不同区域和社群的差异。

账本条目\texttt{TK-0221-SHU-HAN-FOUNDATION}记载，章武元年四月（221年），蜀汉发生“刘备在成都即皇帝位，建元章武，以汉室继承者自居”。这一节点可放在曹丕代汉后，刘备须将汉中王政权转为能够动员益州官民的皇帝体制，并回应关羽败亡后的合法性危机的条件下理解；其后续影响被账本概括为成都朝廷形成与魏相竞争的汉统主张；其可用资源仍主要限于益州和汉中，称帝并未恢复荆州。就地方尺度而言，事件并不只属于朝廷或统帅：征发、道路、仓储、户籍、市场、寺院与家庭关系都可能决定命令如何抵达、战事如何延长，或接收何以在不同地区呈现不一样的速度。

本条列举的核心材料为《三国志·先主传》；《三国志·后主传》；《华阳国志·刘先主志》。这些材料较适合钉住事件的发生、官方表述或特定人物、地点，却不自动构成对全境人口、收入、兵数、信仰或动机的调查。账本的置信与材料提示是“即位年月可靠；国号称谓、礼仪细节及其对各地实际承认的效力有争议”；来源限度进一步说明“蜀汉建国、正统表述与益州政权整合研究”。因此，不能把条目中的政权名称、行政分类或战报修辞直接转译为固定族群和统一社会意志，也不能将制度宣布写成各地同步实施。

在叙事上，更稳妥的做法是把这一事件视为一条需要地方中介才能运作的链：中央的命令、地方官与精英的选择、运输与劳力的条件、以及普通家庭可能面对的迁徙、赋役或安全风险彼此相连。现有材料能够显示其中若干环节，却保留了大量沉默。承认这种沉默并非放弃解释，而是避免以一条编年记录覆盖不同区域和社群的差异。

账本条目\texttt{TK-0221-LIU-BEI-EASTWARD}记载，章武元年（221年），蜀汉与孙吴发生“刘备集结益州军东下，准备以夺回荆州和为关羽报仇为由进攻孙吴”。这一节点可放在孙权夺取荆州并杀关羽，既切断蜀汉东部通道，也使刘备面临集团内部的复仇和领土诉求的条件下理解；其后续影响被账本概括为蜀军离开益州山地防线进入长江中游，孙刘共同抗魏的联盟关系被战争取代。就地方尺度而言，事件并不只属于朝廷或统帅：征发、道路、仓储、户籍、市场、寺院与家庭关系都可能决定命令如何抵达、战事如何延长，或接收何以在不同地区呈现不一样的速度。

本条列举的核心材料为《三国志·先主传》；《三国志·张飞传》；《三国志·陆逊传》。这些材料较适合钉住事件的发生、官方表述或特定人物、地点，却不自动构成对全境人口、收入、兵数、信仰或动机的调查。账本的置信与材料提示是“出兵方向可靠；兵力、动机轻重及张飞遇害对战局的影响有争议”；来源限度进一步说明“夷陵战前动员、荆州归属与蜀汉军事决策研究”。因此，不能把条目中的政权名称、行政分类或战报修辞直接转译为固定族群和统一社会意志，也不能将制度宣布写成各地同步实施。

在叙事上，更稳妥的做法是把这一事件视为一条需要地方中介才能运作的链：中央的命令、地方官与精英的选择、运输与劳力的条件、以及普通家庭可能面对的迁徙、赋役或安全风险彼此相连。现有材料能够显示其中若干环节，却保留了大量沉默。承认这种沉默并非放弃解释，而是避免以一条编年记录覆盖不同区域和社群的差异。

账本条目\texttt{TK-0222-WEI-INVESTS-SUN-QUAN}记载，黄初三年（222年），曹魏与孙权集团发生“曹丕册封孙权为吴王，试图以册命、质子要求和荆州牧名义约束长江下游政权”。这一节点可放在魏在刘备东征时希望利用孙刘冲突，以名义封授换取孙权承认北方皇帝和交出人质的条件下理解；其后续影响被账本概括为孙权暂借魏的封号牵制蜀汉而未交出统治权；魏吴间的宗藩语言很快让位于边境战争。就地方尺度而言，事件并不只属于朝廷或统帅：征发、道路、仓储、户籍、市场、寺院与家庭关系都可能决定命令如何抵达、战事如何延长，或接收何以在不同地区呈现不一样的速度。

本条列举的核心材料为《三国志·吴主传》；《三国志·文帝纪》；《资治通鉴》魏纪。这些材料较适合钉住事件的发生、官方表述或特定人物、地点，却不自动构成对全境人口、收入、兵数、信仰或动机的调查。账本的置信与材料提示是“册封事实可靠；质子交涉的先后、封授实际约束力与孙权的受命姿态有争议”；来源限度进一步说明“魏—吴册封关系、质子外交与长江政治研究”。因此，不能把条目中的政权名称、行政分类或战报修辞直接转译为固定族群和统一社会意志，也不能将制度宣布写成各地同步实施。

在叙事上，更稳妥的做法是把这一事件视为一条需要地方中介才能运作的链：中央的命令、地方官与精英的选择、运输与劳力的条件、以及普通家庭可能面对的迁徙、赋役或安全风险彼此相连。现有材料能够显示其中若干环节，却保留了大量沉默。承认这种沉默并非放弃解释，而是避免以一条编年记录覆盖不同区域和社群的差异。

账本条目\texttt{TK-0222-YILING}记载，章武二年（222年），蜀汉与孙吴发生“陆逊统吴军在夷陵、猇亭一线击破刘备，蜀军沿长江退向白帝”。这一节点可放在刘备军深入荆州后补给线拉长且营垒连结，孙权以陆逊统合江东将领并等待蜀军疲敝的条件下理解；其后续影响被账本概括为蜀汉失去夺回荆州的机会并元气受损；孙吴守住上游屏障，魏无法再期待孙刘相互消灭后立即接收南方。就地方尺度而言，事件并不只属于朝廷或统帅：征发、道路、仓储、户籍、市场、寺院与家庭关系都可能决定命令如何抵达、战事如何延长，或接收何以在不同地区呈现不一样的速度。

本条列举的核心材料为《三国志·先主传》；《三国志·陆逊传》；《三国志·吴主传》；《资治通鉴》魏纪。这些材料较适合钉住事件的发生、官方表述或特定人物、地点，却不自动构成对全境人口、收入、兵数、信仰或动机的调查。账本的置信与材料提示是“胜负与大致地点可靠；火攻展开、伤亡、营垒长度和陆逊受权过程有争议”；来源限度进一步说明“夷陵之战、长江峡谷战场与孙刘联盟破裂研究”。因此，不能把条目中的政权名称、行政分类或战报修辞直接转译为固定族群和统一社会意志，也不能将制度宣布写成各地同步实施。

在叙事上，更稳妥的做法是把这一事件视为一条需要地方中介才能运作的链：中央的命令、地方官与精英的选择、运输与劳力的条件、以及普通家庭可能面对的迁徙、赋役或安全风险彼此相连。现有材料能够显示其中若干环节，却保留了大量沉默。承认这种沉默并非放弃解释，而是避免以一条编年记录覆盖不同区域和社群的差异。

账本条目\texttt{TK-0223-LIU-BEI-DEATH}记载，章武三年四月（223年），蜀汉发生“刘备在永安去世，刘禅继位，诸葛亮受遗诏辅政”。这一节点可放在夷陵败退后刘备驻永安，蜀汉既须防孙吴又须避免益州内部在皇帝更替时分裂的条件下理解；其后续影响被账本概括为年幼的刘禅继位而诸葛亮掌理军政，蜀汉将从复仇性东征转向修复联盟、安定南中和守卫汉中。就地方尺度而言，事件并不只属于朝廷或统帅：征发、道路、仓储、户籍、市场、寺院与家庭关系都可能决定命令如何抵达、战事如何延长，或接收何以在不同地区呈现不一样的速度。

本条列举的核心材料为《三国志·先主传》；《三国志·后主传》；《三国志·诸葛亮传》。这些材料较适合钉住事件的发生、官方表述或特定人物、地点，却不自动构成对全境人口、收入、兵数、信仰或动机的调查。账本的置信与材料提示是“卒年、继位和辅政事实可靠；遗诏措辞、永安部署及刘备病因有争议”；来源限度进一步说明“蜀汉继承、永安托孤与丞相辅政研究”。因此，不能把条目中的政权名称、行政分类或战报修辞直接转译为固定族群和统一社会意志，也不能将制度宣布写成各地同步实施。

在叙事上，更稳妥的做法是把这一事件视为一条需要地方中介才能运作的链：中央的命令、地方官与精英的选择、运输与劳力的条件、以及普通家庭可能面对的迁徙、赋役或安全风险彼此相连。现有材料能够显示其中若干环节，却保留了大量沉默。承认这种沉默并非放弃解释，而是避免以一条编年记录覆盖不同区域和社群的差异。

账本条目\texttt{TK-0223-SHU-WU-RECONCILIATION}记载，建兴元年（223年），蜀汉与孙吴发生“邓芝奉使孙吴，孙权重新与蜀汉通好，双方恢复共同抗魏的外交框架”。这一节点可放在夷陵后蜀汉孤立而魏已在淮汉方向施压，孙权也不能接受魏以册封和人质要求持续干预江东的条件下理解；其后续影响被账本概括为长江上游战线暂获稳定，蜀汉可转向南中和汉中；联盟仍未消除荆州归属造成的结构性裂痕。就地方尺度而言，事件并不只属于朝廷或统帅：征发、道路、仓储、户籍、市场、寺院与家庭关系都可能决定命令如何抵达、战事如何延长，或接收何以在不同地区呈现不一样的速度。

本条列举的核心材料为《三国志·邓芝传》；《三国志·吴主传》；《三国志·诸葛亮传》。这些材料较适合钉住事件的发生、官方表述或特定人物、地点，却不自动构成对全境人口、收入、兵数、信仰或动机的调查。账本的置信与材料提示是“遣使和通好结果可靠；盟约条款、使节言辞及实际协同程度有争议”；来源限度进一步说明“邓芝使吴、孙刘复盟与三国外交研究”。因此，不能把条目中的政权名称、行政分类或战报修辞直接转译为固定族群和统一社会意志，也不能将制度宣布写成各地同步实施。

在叙事上，更稳妥的做法是把这一事件视为一条需要地方中介才能运作的链：中央的命令、地方官与精英的选择、运输与劳力的条件、以及普通家庭可能面对的迁徙、赋役或安全风险彼此相连。现有材料能够显示其中若干环节，却保留了大量沉默。承认这种沉默并非放弃解释，而是避免以一条编年记录覆盖不同区域和社群的差异。

账本条目\texttt{TK-0224-CAO-PI-GUANGLING}记载，黄初五年（224年），曹魏与孙吴发生“曹丕至广陵观察南征，因长江冰冻、舟师与渡江条件不具而撤军”。这一节点可放在魏欲趁孙吴与蜀汉刚恢复关系的脆弱阶段试探渡江，却缺少熟悉江面的水军和稳定前进基地的条件下理解；其后续影响被账本概括为长江水文和舟师能力限制魏的南下计划，三国对峙不仅取决于兵额，也受季节和交通技术制约。就地方尺度而言，事件并不只属于朝廷或统帅：征发、道路、仓储、户籍、市场、寺院与家庭关系都可能决定命令如何抵达、战事如何延长，或接收何以在不同地区呈现不一样的速度。

本条列举的核心材料为《三国志·文帝纪》；《三国志·吴主传》；《资治通鉴》魏纪。这些材料较适合钉住事件的发生、官方表述或特定人物、地点，却不自动构成对全境人口、收入、兵数、信仰或动机的调查。账本的置信与材料提示是“行幸广陵与撤军可靠；冰情、军队规模和曹丕战略意图有争议”；来源限度进一步说明“曹丕广陵南征、长江环境与魏吴边防研究”。因此，不能把条目中的政权名称、行政分类或战报修辞直接转译为固定族群和统一社会意志，也不能将制度宣布写成各地同步实施。

在叙事上，更稳妥的做法是把这一事件视为一条需要地方中介才能运作的链：中央的命令、地方官与精英的选择、运输与劳力的条件、以及普通家庭可能面对的迁徙、赋役或安全风险彼此相连。现有材料能够显示其中若干环节，却保留了大量沉默。承认这种沉默并非放弃解释，而是避免以一条编年记录覆盖不同区域和社群的差异。

账本条目\texttt{TK-0225-SHU-NANZHONG}记载，建兴三年（225年），蜀汉与南中地方集团发生“诸葛亮南征，平定益州郡、永昌等地叛乱并重置地方统治关系”。这一节点可放在刘备死后南中豪强、夷帅与郡县官之间的紧张关系爆发，蜀汉必须恢复铜铁、赋役和通往西南的交通的条件下理解；其后续影响被账本概括为成都重新取得南中资源和人力，但史书中的归附叙事不能等同地方社会被完全整合；孟获故事须分层处理。就地方尺度而言，事件并不只属于朝廷或统帅：征发、道路、仓储、户籍、市场、寺院与家庭关系都可能决定命令如何抵达、战事如何延长，或接收何以在不同地区呈现不一样的速度。

本条列举的核心材料为《三国志·诸葛亮传》；《三国志·马谡传》裴松之注引《襄阳记》；《华阳国志·南中志》。这些材料较适合钉住事件的发生、官方表述或特定人物、地点，却不自动构成对全境人口、收入、兵数、信仰或动机的调查。账本的置信与材料提示是“南征及其大致结果可靠；孟获身份、七擒次数、行军路线和改置措施的细节有争议”；来源限度进一步说明“蜀汉南中经营、地方首领与后出战功叙事研究”。因此，不能把条目中的政权名称、行政分类或战报修辞直接转译为固定族群和统一社会意志，也不能将制度宣布写成各地同步实施。

在叙事上，更稳妥的做法是把这一事件视为一条需要地方中介才能运作的链：中央的命令、地方官与精英的选择、运输与劳力的条件、以及普通家庭可能面对的迁徙、赋役或安全风险彼此相连。现有材料能够显示其中若干环节，却保留了大量沉默。承认这种沉默并非放弃解释，而是避免以一条编年记录覆盖不同区域和社群的差异。

账本条目\texttt{TK-0226-CAO-PI-DEATH}记载，黄初七年五月（226年），曹魏发生“曹丕去世，曹叡继位，是为魏明帝”。这一节点可放在魏的皇位继承已在曹丕在世时确定，但对吴、蜀的边境战争和宗室、外戚、近侍并存使新帝需要迅速确认中枢命令的条件下理解；其后续影响被账本概括为曹叡承接魏的北方国家机器，魏蜀汉中和魏吴淮南战场的竞争将转入新一轮部署。就地方尺度而言，事件并不只属于朝廷或统帅：征发、道路、仓储、户籍、市场、寺院与家庭关系都可能决定命令如何抵达、战事如何延长，或接收何以在不同地区呈现不一样的速度。

本条列举的核心材料为《三国志·文帝纪》；《三国志·明帝纪》；《资治通鉴》魏纪。这些材料较适合钉住事件的发生、官方表述或特定人物、地点，却不自动构成对全境人口、收入、兵数、信仰或动机的调查。账本的置信与材料提示是“卒年与继位可靠；遗诏所定辅政层级和各集团的权力边界有争议”；来源限度进一步说明“魏文帝身后继承、曹叡即位与魏初政治研究”。因此，不能把条目中的政权名称、行政分类或战报修辞直接转译为固定族群和统一社会意志，也不能将制度宣布写成各地同步实施。

在叙事上，更稳妥的做法是把这一事件视为一条需要地方中介才能运作的链：中央的命令、地方官与精英的选择、运输与劳力的条件、以及普通家庭可能面对的迁徙、赋役或安全风险彼此相连。现有材料能够显示其中若干环节，却保留了大量沉默。承认这种沉默并非放弃解释，而是避免以一条编年记录覆盖不同区域和社群的差异。

账本条目\texttt{TK-0227-ZHUGE-LIANG-HANZHONG}记载，建兴五年（227年），蜀汉发生“诸葛亮驻汉中并上《出师表》，以北伐准备重新组织蜀汉军政”。这一节点可放在南中大体安定、孙吴复盟后，蜀汉希望利用汉中山口牵制魏，并以进取战略维系小国的军政凝聚的条件下理解；其后续影响被账本概括为蜀汉将资源集中于陇右方向；北伐的补给距离和关隘限制随即成为蜀汉国家能力的持续约束。就地方尺度而言，事件并不只属于朝廷或统帅：征发、道路、仓储、户籍、市场、寺院与家庭关系都可能决定命令如何抵达、战事如何延长，或接收何以在不同地区呈现不一样的速度。

本条列举的核心材料为《三国志·诸葛亮传》；《三国志·后主传》；《出师表》传本。这些材料较适合钉住事件的发生、官方表述或特定人物、地点，却不自动构成对全境人口、收入、兵数、信仰或动机的调查。账本的置信与材料提示是“驻汉中和北伐准备可靠；《出师表》成文语境、传本与其是否完整代表当时决策讨论有争议”；来源限度进一步说明“出师表、汉中驻军与蜀汉北伐动员研究”。因此，不能把条目中的政权名称、行政分类或战报修辞直接转译为固定族群和统一社会意志，也不能将制度宣布写成各地同步实施。

在叙事上，更稳妥的做法是把这一事件视为一条需要地方中介才能运作的链：中央的命令、地方官与精英的选择、运输与劳力的条件、以及普通家庭可能面对的迁徙、赋役或安全风险彼此相连。现有材料能够显示其中若干环节，却保留了大量沉默。承认这种沉默并非放弃解释，而是避免以一条编年记录覆盖不同区域和社群的差异。

\subsection{道路、边地与人口移动}

账本条目\texttt{TK-0228-FIRST-NORTHERN-EXPEDITION}记载，建兴六年春（228年），蜀汉与曹魏发生“诸葛亮出祁山，南安、天水、安定一度响应蜀汉，魏调张郃等西援”。这一节点可放在蜀汉试图趁魏对陇右防备不足，从汉中越山进入渭水上游并争取地方官民的条件下理解；其后续影响被账本概括为魏以中央调兵和陇右守将稳住局面；短暂响应显示边郡政治并非铁板一块，却未转化为蜀汉常态控制。就地方尺度而言，事件并不只属于朝廷或统帅：征发、道路、仓储、户籍、市场、寺院与家庭关系都可能决定命令如何抵达、战事如何延长，或接收何以在不同地区呈现不一样的速度。

本条列举的核心材料为《三国志·诸葛亮传》；《三国志·明帝纪》；《三国志·张郃传》；《资治通鉴》魏纪。这些材料较适合钉住事件的发生、官方表述或特定人物、地点，却不自动构成对全境人口、收入、兵数、信仰或动机的调查。账本的置信与材料提示是“出兵与数郡震动可靠；响应范围、地方倒向的原因和双方兵力有争议”；来源限度进一步说明“第一次北伐、陇右郡县与蜀魏动员研究”。因此，不能把条目中的政权名称、行政分类或战报修辞直接转译为固定族群和统一社会意志，也不能将制度宣布写成各地同步实施。

在叙事上，更稳妥的做法是把这一事件视为一条需要地方中介才能运作的链：中央的命令、地方官与精英的选择、运输与劳力的条件、以及普通家庭可能面对的迁徙、赋役或安全风险彼此相连。现有材料能够显示其中若干环节，却保留了大量沉默。承认这种沉默并非放弃解释，而是避免以一条编年记录覆盖不同区域和社群的差异。

账本条目\texttt{TK-0228-JIETING}记载，建兴六年（228年），蜀汉与曹魏发生“马谡失守街亭，诸葛亮撤回汉中并处置失军者”。这一节点可放在祁山方向的蜀军须守住通往陇右的要道，马谡部与魏援军争夺水源和高地时失去稳定阵地的条件下理解；其后续影响被账本概括为第一次北伐无法维持，蜀汉重新收缩至汉中；后世将失败归为单一将领失误的叙事须与补给和指挥结构并读。就地方尺度而言，事件并不只属于朝廷或统帅：征发、道路、仓储、户籍、市场、寺院与家庭关系都可能决定命令如何抵达、战事如何延长，或接收何以在不同地区呈现不一样的速度。

本条列举的核心材料为《三国志·马谡传》；《三国志·王平传》；《三国志·张郃传》；《三国志·诸葛亮传》裴松之注。这些材料较适合钉住事件的发生、官方表述或特定人物、地点，却不自动构成对全境人口、收入、兵数、信仰或动机的调查。账本的置信与材料提示是“失守与撤军可靠；马谡部署、王平作用、军法程序和责任归属有争议”；来源限度进一步说明“街亭失守、陇右补给线与蜀汉军纪研究”。因此，不能把条目中的政权名称、行政分类或战报修辞直接转译为固定族群和统一社会意志，也不能将制度宣布写成各地同步实施。

在叙事上，更稳妥的做法是把这一事件视为一条需要地方中介才能运作的链：中央的命令、地方官与精英的选择、运输与劳力的条件、以及普通家庭可能面对的迁徙、赋役或安全风险彼此相连。现有材料能够显示其中若干环节，却保留了大量沉默。承认这种沉默并非放弃解释，而是避免以一条编年记录覆盖不同区域和社群的差异。

账本条目\texttt{TK-0228-MENG-DA-REBELLION}记载，太和二年（228年），曹魏发生“新城太守孟达谋反，司马懿急行军围新城，孟达为部下所杀”。这一节点可放在孟达处在魏蜀交通边缘，对魏的任用与自身安全缺乏信任，蜀汉北伐又改变了其对形势的判断的条件下理解；其后续影响被账本概括为魏以快速机动保住汉水上游节点，司马懿在西部军政中的声望上升；相关书信不能未经辨伪视为直接内情。就地方尺度而言，事件并不只属于朝廷或统帅：征发、道路、仓储、户籍、市场、寺院与家庭关系都可能决定命令如何抵达、战事如何延长，或接收何以在不同地区呈现不一样的速度。

本条列举的核心材料为《三国志·明帝纪》；《三国志·司马懿传》；《三国志·孟达传》裴松之注引《魏略》。这些材料较适合钉住事件的发生、官方表述或特定人物、地点，却不自动构成对全境人口、收入、兵数、信仰或动机的调查。账本的置信与材料提示是“叛乱、围城和孟达死亡可靠；书信真伪、行军日程和孟达与诸葛亮联络程度有争议”；来源限度进一步说明“孟达叛乱、上庸新城与魏蜀边境情报研究”。因此，不能把条目中的政权名称、行政分类或战报修辞直接转译为固定族群和统一社会意志，也不能将制度宣布写成各地同步实施。

在叙事上，更稳妥的做法是把这一事件视为一条需要地方中介才能运作的链：中央的命令、地方官与精英的选择、运输与劳力的条件、以及普通家庭可能面对的迁徙、赋役或安全风险彼此相连。现有材料能够显示其中若干环节，却保留了大量沉默。承认这种沉默并非放弃解释，而是避免以一条编年记录覆盖不同区域和社群的差异。

账本条目\texttt{TK-0228-SHITING}记载，太和二年（228年），孙吴与曹魏发生“陆逊在石亭击败曹休所率魏军，魏的江淮攻势受挫”。这一节点可放在魏试图沿江淮推进并利用吴方内部关系，孙吴以陆逊、朱桓等将领掌握水陆节点和情报优势的条件下理解；其后续影响被账本概括为魏军后撤并调整淮南防御，孙吴证明其可在江北作战；诈降情节主要见于传记叙事，须区分战果与戏剧化细节。就地方尺度而言，事件并不只属于朝廷或统帅：征发、道路、仓储、户籍、市场、寺院与家庭关系都可能决定命令如何抵达、战事如何延长，或接收何以在不同地区呈现不一样的速度。

本条列举的核心材料为《三国志·吴主传》；《三国志·陆逊传》；《三国志·曹休传》；《资治通鉴》魏纪。这些材料较适合钉住事件的发生、官方表述或特定人物、地点，却不自动构成对全境人口、收入、兵数、信仰或动机的调查。账本的置信与材料提示是“战役结果可靠；诈降信息、参战兵力、战场位置和追击范围有争议”；来源限度进一步说明“石亭之战、江淮情报与魏吴战争研究”。因此，不能把条目中的政权名称、行政分类或战报修辞直接转译为固定族群和统一社会意志，也不能将制度宣布写成各地同步实施。

在叙事上，更稳妥的做法是把这一事件视为一条需要地方中介才能运作的链：中央的命令、地方官与精英的选择、运输与劳力的条件、以及普通家庭可能面对的迁徙、赋役或安全风险彼此相连。现有材料能够显示其中若干环节，却保留了大量沉默。承认这种沉默并非放弃解释，而是避免以一条编年记录覆盖不同区域和社群的差异。

账本条目\texttt{TK-0229-WUDU-YINPING}记载，建兴七年（229年），蜀汉与曹魏发生“诸葛亮取武都、阴平二郡，魏将郭淮退却，蜀汉扩展汉中西北屏障”。这一节点可放在蜀汉第一次北伐失败后仍须改善汉中西北的关隘防御，并寻找较短的局部进攻机会的条件下理解；其后续影响被账本概括为蜀汉取得两郡名义和部分山地通道，但难以由此突破魏的渭水防线；边郡控制范围须逐县核验。就地方尺度而言，事件并不只属于朝廷或统帅：征发、道路、仓储、户籍、市场、寺院与家庭关系都可能决定命令如何抵达、战事如何延长，或接收何以在不同地区呈现不一样的速度。

本条列举的核心材料为《三国志·诸葛亮传》；《三国志·郭淮传》；《资治通鉴》魏纪。这些材料较适合钉住事件的发生、官方表述或特定人物、地点，却不自动构成对全境人口、收入、兵数、信仰或动机的调查。账本的置信与材料提示是“取郡结果大体可靠；郡界、守令转换和郭淮撤军原因有争议”；来源限度进一步说明“武都阴平争夺、陇西山地与蜀汉边防研究”。因此，不能把条目中的政权名称、行政分类或战报修辞直接转译为固定族群和统一社会意志，也不能将制度宣布写成各地同步实施。

在叙事上，更稳妥的做法是把这一事件视为一条需要地方中介才能运作的链：中央的命令、地方官与精英的选择、运输与劳力的条件、以及普通家庭可能面对的迁徙、赋役或安全风险彼此相连。现有材料能够显示其中若干环节，却保留了大量沉默。承认这种沉默并非放弃解释，而是避免以一条编年记录覆盖不同区域和社群的差异。

账本条目\texttt{TK-0229-SUN-QUAN-EMPEROR}记载，黄龙元年四月（229年），孙吴发生“孙权在武昌即皇帝位，建元黄龙，孙吴正式提出独立帝国名号”。这一节点可放在魏的册封与质子要求不能满足孙权的自主统治，蜀汉亦已称帝，江东政权需要以独立名号配置官爵和外交的条件下理解；其后续影响被账本概括为孙吴与魏、蜀汉形成三个皇帝中心；称帝强化建国仪式，却未消除长江中上游与淮南的军事压力。就地方尺度而言，事件并不只属于朝廷或统帅：征发、道路、仓储、户籍、市场、寺院与家庭关系都可能决定命令如何抵达、战事如何延长，或接收何以在不同地区呈现不一样的速度。

本条列举的核心材料为《三国志·吴主传》；《三国志·孙登传》；《资治通鉴》魏纪。这些材料较适合钉住事件的发生、官方表述或特定人物、地点，却不自动构成对全境人口、收入、兵数、信仰或动机的调查。账本的置信与材料提示是“即位年月可靠；受禅式礼仪、群臣参与和对地方的实际动员力有争议”；来源限度进一步说明“孙吴称帝、黄龙年号与江东合法性研究”。因此，不能把条目中的政权名称、行政分类或战报修辞直接转译为固定族群和统一社会意志，也不能将制度宣布写成各地同步实施。

在叙事上，更稳妥的做法是把这一事件视为一条需要地方中介才能运作的链：中央的命令、地方官与精英的选择、运输与劳力的条件、以及普通家庭可能面对的迁徙、赋役或安全风险彼此相连。现有材料能够显示其中若干环节，却保留了大量沉默。承认这种沉默并非放弃解释，而是避免以一条编年记录覆盖不同区域和社群的差异。

账本条目\texttt{TK-0229-JIANYE-CAPITAL}记载，黄龙元年（229年），孙吴发生“孙权将都城由武昌迁回建业，长江下游成为孙吴中枢”。这一节点可放在孙吴称帝后须在江东腹地建立较稳定的官僚与交通中心，同时保留武昌对上游的军事作用的条件下理解；其后续影响被账本概括为建业成为孙吴及后世建康政治传统的起点；城址材料可校验城市演变，不能直接复原全部三国宫城布局。就地方尺度而言，事件并不只属于朝廷或统帅：征发、道路、仓储、户籍、市场、寺院与家庭关系都可能决定命令如何抵达、战事如何延长，或接收何以在不同地区呈现不一样的速度。

本条列举的核心材料为《三国志·吴主传》；《建康实录》；南京地区六朝城址考古材料。这些材料较适合钉住事件的发生、官方表述或特定人物、地点，却不自动构成对全境人口、收入、兵数、信仰或动机的调查。账本的置信与材料提示是“迁都可靠；宫城建设、迁徙人口与武昌—建业双中心的实际分工有争议”；来源限度进一步说明“孙吴建业建都、长江下游城市与区域资源研究”。因此，不能把条目中的政权名称、行政分类或战报修辞直接转译为固定族群和统一社会意志，也不能将制度宣布写成各地同步实施。

在叙事上，更稳妥的做法是把这一事件视为一条需要地方中介才能运作的链：中央的命令、地方官与精英的选择、运输与劳力的条件、以及普通家庭可能面对的迁徙、赋役或安全风险彼此相连。现有材料能够显示其中若干环节，却保留了大量沉默。承认这种沉默并非放弃解释，而是避免以一条编年记录覆盖不同区域和社群的差异。

账本条目\texttt{TK-0230-WEI-WEN-MARITIME}记载，黄龙二年（230年），孙吴与海上岛屿社会发生“孙权遣卫温、诸葛直率舟师求夷洲、亶洲，航海军民多病死而未获预期成果”。这一节点可放在孙吴拥有沿海舟师并寻求扩展人力、资源和外部信息，但对外海环境与岛屿社会缺乏稳定补给条件的条件下理解；其后续影响被账本概括为远航显示孙吴政权的海上动员能力及其限度；夷洲是否对应今台湾等问题必须保留文本地理的不确定性。就地方尺度而言，事件并不只属于朝廷或统帅：征发、道路、仓储、户籍、市场、寺院与家庭关系都可能决定命令如何抵达、战事如何延长，或接收何以在不同地区呈现不一样的速度。

本条列举的核心材料为《三国志·吴主传》；《三国志·吴主传》裴松之注引《临海水土志》；《资治通鉴》魏纪。这些材料较适合钉住事件的发生、官方表述或特定人物、地点，却不自动构成对全境人口、收入、兵数、信仰或动机的调查。账本的置信与材料提示是“遣航与损失大体可靠；夷洲、亶洲的地望、航程和人数不可据后世地名直接断定”；来源限度进一步说明“孙吴海上远航、夷洲亶洲地望与海洋史研究”。因此，不能把条目中的政权名称、行政分类或战报修辞直接转译为固定族群和统一社会意志，也不能将制度宣布写成各地同步实施。

在叙事上，更稳妥的做法是把这一事件视为一条需要地方中介才能运作的链：中央的命令、地方官与精英的选择、运输与劳力的条件、以及普通家庭可能面对的迁徙、赋役或安全风险彼此相连。现有材料能够显示其中若干环节，却保留了大量沉默。承认这种沉默并非放弃解释，而是避免以一条编年记录覆盖不同区域和社群的差异。

账本条目\texttt{TK-0230-CAO-ZHEN-HANZHONG}记载，太和四年至五年（230年至231年），曹魏与蜀汉发生“曹真等自关中多路准备攻蜀，栈道与山道遭大雨阻断，魏军未能深入汉中”。这一节点可放在魏希望利用兵力与关中资源压迫蜀汉，却必须跨越秦岭并维持长距离栈道运输的条件下理解；其后续影响被账本概括为天气和道路使魏的大规模攻势无法兑现，汉中山地继续提供蜀汉防御纵深；传记中的工程成效须与地形分别评估。就地方尺度而言，事件并不只属于朝廷或统帅：征发、道路、仓储、户籍、市场、寺院与家庭关系都可能决定命令如何抵达、战事如何延长，或接收何以在不同地区呈现不一样的速度。

本条列举的核心材料为《三国志·明帝纪》；《三国志·曹真传》；《三国志·王昶传》；《资治通鉴》魏纪。这些材料较适合钉住事件的发生、官方表述或特定人物、地点，却不自动构成对全境人口、收入、兵数、信仰或动机的调查。账本的置信与材料提示是“出师与因雨撤军可靠；各路行军、道路工程和损失规模有争议”；来源限度进一步说明“曹真伐蜀、褒斜山道与山地后勤研究”。因此，不能把条目中的政权名称、行政分类或战报修辞直接转译为固定族群和统一社会意志，也不能将制度宣布写成各地同步实施。

在叙事上，更稳妥的做法是把这一事件视为一条需要地方中介才能运作的链：中央的命令、地方官与精英的选择、运输与劳力的条件、以及普通家庭可能面对的迁徙、赋役或安全风险彼此相连。现有材料能够显示其中若干环节，却保留了大量沉默。承认这种沉默并非放弃解释，而是避免以一条编年记录覆盖不同区域和社群的差异。

账本条目\texttt{TK-0231-QISHAN-EXPEDITION}记载，建兴九年（231年），蜀汉与曹魏发生“诸葛亮再次出祁山，与司马懿部相持，因粮运不继而撤军”。这一节点可放在蜀汉需持续在陇右制造压力以维持边境主动，但益州至前线的粮道难以支撑久战的条件下理解；其后续影响被账本概括为魏未能消灭蜀军，蜀汉也未能固守陇右；战场结果再次把小国战略的关键置于运输而非单一会战。就地方尺度而言，事件并不只属于朝廷或统帅：征发、道路、仓储、户籍、市场、寺院与家庭关系都可能决定命令如何抵达、战事如何延长，或接收何以在不同地区呈现不一样的速度。

本条列举的核心材料为《三国志·诸葛亮传》；《三国志·司马懿传》；《三国志·张郃传》；《资治通鉴》魏纪。这些材料较适合钉住事件的发生、官方表述或特定人物、地点，却不自动构成对全境人口、收入、兵数、信仰或动机的调查。账本的置信与材料提示是“出兵、对峙和撤军可靠；木牛流马、卤城战果和双方伤亡有争议”；来源限度进一步说明“祁山相持、蜀汉粮运与魏蜀战略研究”。因此，不能把条目中的政权名称、行政分类或战报修辞直接转译为固定族群和统一社会意志，也不能将制度宣布写成各地同步实施。

在叙事上，更稳妥的做法是把这一事件视为一条需要地方中介才能运作的链：中央的命令、地方官与精英的选择、运输与劳力的条件、以及普通家庭可能面对的迁徙、赋役或安全风险彼此相连。现有材料能够显示其中若干环节，却保留了大量沉默。承认这种沉默并非放弃解释，而是避免以一条编年记录覆盖不同区域和社群的差异。

账本条目\texttt{TK-0231-LI-YAN-DISMISSAL}记载，建兴九年（231年），蜀汉发生“李严因粮运失职并伪称奉诏召军，被诸葛亮弹劾废为平民”。这一节点可放在祁山前线依赖永安—益州方面的粮运，撤军使统帅与负责转输的重臣围绕责任和声誉发生冲突的条件下理解；其后续影响被账本概括为蜀汉失去一名资深大臣并强化诸葛亮对军政的控制；案件材料多出自胜方文书，不能把弹劾词当作无争议事实。就地方尺度而言，事件并不只属于朝廷或统帅：征发、道路、仓储、户籍、市场、寺院与家庭关系都可能决定命令如何抵达、战事如何延长，或接收何以在不同地区呈现不一样的速度。

本条列举的核心材料为《三国志·李严传》；《三国志·诸葛亮传》；《三国志·李严传》裴松之注引《诸葛亮集》。这些材料较适合钉住事件的发生、官方表述或特定人物、地点，却不自动构成对全境人口、收入、兵数、信仰或动机的调查。账本的置信与材料提示是“罢黜事实可靠；李严失职、诏令真伪和诸葛亮处置的程序与责任比例有争议”；来源限度进一步说明“蜀汉后勤官僚、李严案与北伐责任研究”。因此，不能把条目中的政权名称、行政分类或战报修辞直接转译为固定族群和统一社会意志，也不能将制度宣布写成各地同步实施。

在叙事上，更稳妥的做法是把这一事件视为一条需要地方中介才能运作的链：中央的命令、地方官与精英的选择、运输与劳力的条件、以及普通家庭可能面对的迁徙、赋役或安全风险彼此相连。现有材料能够显示其中若干环节，却保留了大量沉默。承认这种沉默并非放弃解释，而是避免以一条编年记录覆盖不同区域和社群的差异。

\subsection{社会关系与行政分类}

账本条目\texttt{TK-0232-GONGSUN-YUAN-WU-TIES}记载，太和六年（232年），辽东公孙渊与孙吴发生“公孙渊遣使通孙吴，孙权册拜其为燕王，试图在魏的东北边缘建立外部盟点”。这一节点可放在公孙氏长期据辽东并在魏、吴之间权衡，孙吴则希望以海路外交牵制魏而不直接承担辽东防务的条件下理解；其后续影响被账本概括为吴辽联系短暂打开魏后方的外交风险，也暴露跨海盟约缺乏陆上支援的脆弱性。就地方尺度而言，事件并不只属于朝廷或统帅：征发、道路、仓储、户籍、市场、寺院与家庭关系都可能决定命令如何抵达、战事如何延长，或接收何以在不同地区呈现不一样的速度。

本条列举的核心材料为《三国志·吴主传》；《三国志·公孙度传》；《三国志·公孙度传》裴松之注引《魏略》。这些材料较适合钉住事件的发生、官方表述或特定人物、地点，却不自动构成对全境人口、收入、兵数、信仰或动机的调查。账本的置信与材料提示是“互遣使节和册拜大体可靠；公孙渊真实意图、封号授受程序和辽东控制范围有争议”；来源限度进一步说明“辽东公孙氏、孙吴外交与魏东北边防研究”。因此，不能把条目中的政权名称、行政分类或战报修辞直接转译为固定族群和统一社会意志，也不能将制度宣布写成各地同步实施。

在叙事上，更稳妥的做法是把这一事件视为一条需要地方中介才能运作的链：中央的命令、地方官与精英的选择、运输与劳力的条件、以及普通家庭可能面对的迁徙、赋役或安全风险彼此相连。现有材料能够显示其中若干环节，却保留了大量沉默。承认这种沉默并非放弃解释，而是避免以一条编年记录覆盖不同区域和社群的差异。

账本条目\texttt{TK-0233-GONGSUN-YUAN-ENVOYS}记载，嘉禾二年（233年），辽东公孙渊与孙吴发生“公孙渊杀孙吴使者张弥、许晏并夺取所赐财物，转而向魏请和”。这一节点可放在公孙渊面对魏的近距离军事威胁，判断孙吴不能跨海提供持续援军，遂以处死使者表明立场的条件下理解；其后续影响被账本概括为孙吴失去在辽东的盟点，魏对公孙氏的猜忌未除；这一事件说明使节册封不等同于可执行的军事同盟。就地方尺度而言，事件并不只属于朝廷或统帅：征发、道路、仓储、户籍、市场、寺院与家庭关系都可能决定命令如何抵达、战事如何延长，或接收何以在不同地区呈现不一样的速度。

本条列举的核心材料为《三国志·吴主传》；《三国志·公孙度传》；《三国志·公孙度传》裴松之注引《吴书》。这些材料较适合钉住事件的发生、官方表述或特定人物、地点，却不自动构成对全境人口、收入、兵数、信仰或动机的调查。账本的置信与材料提示是“使者遇害和外交转向可靠；公孙渊决策动机、财物数额和魏方承诺有争议”；来源限度进一步说明“辽东外交反复、吴使遇害与海路联盟研究”。因此，不能把条目中的政权名称、行政分类或战报修辞直接转译为固定族群和统一社会意志，也不能将制度宣布写成各地同步实施。

在叙事上，更稳妥的做法是把这一事件视为一条需要地方中介才能运作的链：中央的命令、地方官与精英的选择、运输与劳力的条件、以及普通家庭可能面对的迁徙、赋役或安全风险彼此相连。现有材料能够显示其中若干环节，却保留了大量沉默。承认这种沉默并非放弃解释，而是避免以一条编年记录覆盖不同区域和社群的差异。

账本条目\texttt{TK-0234-WUZHANG}记载，建兴十二年（234年），蜀汉与曹魏发生“诸葛亮出斜谷屯五丈原，与司马懿相持后病逝，蜀军撤回汉中”。这一节点可放在蜀汉仍试图由关中方向牵制魏，但补给距离和魏军防守使战事难以取得决定性突破的条件下理解；其后续影响被账本概括为诸葛亮去世结束蜀汉以丞相统合北伐的阶段，魏守住关中；后出轶事应与可核实的撤军序列分开。就地方尺度而言，事件并不只属于朝廷或统帅：征发、道路、仓储、户籍、市场、寺院与家庭关系都可能决定命令如何抵达、战事如何延长，或接收何以在不同地区呈现不一样的速度。

本条列举的核心材料为《三国志·诸葛亮传》；《三国志·司马懿传》；《三国志·杨仪传》裴松之注引《魏氏春秋》。这些材料较适合钉住事件的发生、官方表述或特定人物、地点，却不自动构成对全境人口、收入、兵数、信仰或动机的调查。账本的置信与材料提示是“出兵、相持、死亡和撤军可靠；病因、交战规模、司马懿避战动机和“死诸葛走生仲达”叙事有争议”；来源限度进一步说明“五丈原、诸葛亮之死与蜀汉战略限度研究”。因此，不能把条目中的政权名称、行政分类或战报修辞直接转译为固定族群和统一社会意志，也不能将制度宣布写成各地同步实施。

在叙事上，更稳妥的做法是把这一事件视为一条需要地方中介才能运作的链：中央的命令、地方官与精英的选择、运输与劳力的条件、以及普通家庭可能面对的迁徙、赋役或安全风险彼此相连。现有材料能够显示其中若干环节，却保留了大量沉默。承认这种沉默并非放弃解释，而是避免以一条编年记录覆盖不同区域和社群的差异。

账本条目\texttt{TK-0234-WU-HEFEI-ASSAULT}记载，嘉禾三年（234年），孙吴与曹魏发生“孙权亲攻合肥新城，魏守军与援军固守，吴军撤退”。这一节点可放在蜀汉北伐迫使魏分兵关中，孙吴希望在淮南夺取可渡淮的据点并形成协同压力的条件下理解；其后续影响被账本概括为魏守住合肥防线，吴蜀并攻未能改变三国边界；两线同时作战的协调程度不可由后世战略推演假定。就地方尺度而言，事件并不只属于朝廷或统帅：征发、道路、仓储、户籍、市场、寺院与家庭关系都可能决定命令如何抵达、战事如何延长，或接收何以在不同地区呈现不一样的速度。

本条列举的核心材料为《三国志·吴主传》；《三国志·满宠传》；《三国志·明帝纪》；《资治通鉴》魏纪。这些材料较适合钉住事件的发生、官方表述或特定人物、地点，却不自动构成对全境人口、收入、兵数、信仰或动机的调查。账本的置信与材料提示是“攻城与撤军可靠；孙权亲临位置、围城时长、疾病影响和魏援军作用有争议”；来源限度进一步说明“合肥新城攻防、淮南城防与魏吴战争研究”。因此，不能把条目中的政权名称、行政分类或战报修辞直接转译为固定族群和统一社会意志，也不能将制度宣布写成各地同步实施。

在叙事上，更稳妥的做法是把这一事件视为一条需要地方中介才能运作的链：中央的命令、地方官与精英的选择、运输与劳力的条件、以及普通家庭可能面对的迁徙、赋役或安全风险彼此相连。现有材料能够显示其中若干环节，却保留了大量沉默。承认这种沉默并非放弃解释，而是避免以一条编年记录覆盖不同区域和社群的差异。

账本条目\texttt{TK-0234-JIANG-WAN-REGENT}记载，建兴十二年（234年），蜀汉发生“诸葛亮死后，蒋琬录尚书事，继掌蜀汉中枢政务”。这一节点可放在诸葛亮长期集中军政，去世后蜀汉必须在不引发益州官僚分裂的条件下重组决策和前线防务的条件下理解；其后续影响被账本概括为蒋琬以尚书机构延续辅政，蜀汉北伐节奏趋于审慎；政权并非因一位权臣之死立即失去全部行政能力。就地方尺度而言，事件并不只属于朝廷或统帅：征发、道路、仓储、户籍、市场、寺院与家庭关系都可能决定命令如何抵达、战事如何延长，或接收何以在不同地区呈现不一样的速度。

本条列举的核心材料为《三国志·蒋琬传》；《三国志·后主传》；《华阳国志·刘后主志》。这些材料较适合钉住事件的发生、官方表述或特定人物、地点，却不自动构成对全境人口、收入、兵数、信仰或动机的调查。账本的置信与材料提示是“接掌中枢事实可靠；蒋琬权力起点、与费祎和姜维的分工及政策连续性有争议”；来源限度进一步说明“蒋琬辅政、蜀汉继承与战后行政研究”。因此，不能把条目中的政权名称、行政分类或战报修辞直接转译为固定族群和统一社会意志，也不能将制度宣布写成各地同步实施。

在叙事上，更稳妥的做法是把这一事件视为一条需要地方中介才能运作的链：中央的命令、地方官与精英的选择、运输与劳力的条件、以及普通家庭可能面对的迁徙、赋役或安全风险彼此相连。现有材料能够显示其中若干环节，却保留了大量沉默。承认这种沉默并非放弃解释，而是避免以一条编年记录覆盖不同区域和社群的差异。

账本条目\texttt{TK-0237-GONGSUN-YUAN-REBELLS}记载，景初元年（237年），曹魏与辽东公孙渊发生“公孙渊拒绝入朝受命，自立燕王并与魏公开决裂”。这一节点可放在魏要求公孙渊入朝并重整辽东任命，触及公孙氏长期积累的地方军政自主权的条件下理解；其后续影响被账本概括为辽东危机由暧昧附属关系转为战争，魏须在关中、淮南之外投入资源以解决东北割据中心。就地方尺度而言，事件并不只属于朝廷或统帅：征发、道路、仓储、户籍、市场、寺院与家庭关系都可能决定命令如何抵达、战事如何延长，或接收何以在不同地区呈现不一样的速度。

本条列举的核心材料为《三国志·明帝纪》；《三国志·公孙度传》；《三国志·公孙度传》裴松之注引《魏略》。这些材料较适合钉住事件的发生、官方表述或特定人物、地点，却不自动构成对全境人口、收入、兵数、信仰或动机的调查。账本的置信与材料提示是“反魏与自立名号可靠；称王时间、官号体系和辽东各地服从程度有争议”；来源限度进一步说明“公孙渊叛魏、辽东政权与东北交通研究”。因此，不能把条目中的政权名称、行政分类或战报修辞直接转译为固定族群和统一社会意志，也不能将制度宣布写成各地同步实施。

在叙事上，更稳妥的做法是把这一事件视为一条需要地方中介才能运作的链：中央的命令、地方官与精英的选择、运输与劳力的条件、以及普通家庭可能面对的迁徙、赋役或安全风险彼此相连。现有材料能够显示其中若干环节，却保留了大量沉默。承认这种沉默并非放弃解释，而是避免以一条编年记录覆盖不同区域和社群的差异。

账本条目\texttt{TK-0238-SIMA-YI-LIAODONG}记载，景初二年（238年），曹魏与辽东燕政权发生“司马懿围襄平，公孙渊败死，魏结束公孙氏在辽东的长期割据”。这一节点可放在公孙渊公开称王迫使魏不能再以册封维持边缘秩序，司马懿获授较大军权统率远征的条件下理解；其后续影响被账本概括为魏取得辽东名义控制并消除吴的潜在盟友，但战争暴力和交通距离使战后治理不能简化为立即稳定。就地方尺度而言，事件并不只属于朝廷或统帅：征发、道路、仓储、户籍、市场、寺院与家庭关系都可能决定命令如何抵达、战事如何延长，或接收何以在不同地区呈现不一样的速度。

本条列举的核心材料为《三国志·明帝纪》；《三国志·司马懿传》；《三国志·公孙度传》；《资治通鉴》魏纪。这些材料较适合钉住事件的发生、官方表述或特定人物、地点，却不自动构成对全境人口、收入、兵数、信仰或动机的调查。账本的置信与材料提示是“征讨、围城和公孙渊死亡可靠；雨水、屠杀规模、投降处置和战后行政恢复速度有争议”；来源限度进一步说明“辽东之役、襄平围城与魏东北整合研究”。因此，不能把条目中的政权名称、行政分类或战报修辞直接转译为固定族群和统一社会意志，也不能将制度宣布写成各地同步实施。

在叙事上，更稳妥的做法是把这一事件视为一条需要地方中介才能运作的链：中央的命令、地方官与精英的选择、运输与劳力的条件、以及普通家庭可能面对的迁徙、赋役或安全风险彼此相连。现有材料能够显示其中若干环节，却保留了大量沉默。承认这种沉默并非放弃解释，而是避免以一条编年记录覆盖不同区域和社群的差异。

账本条目\texttt{TK-0238-HIMIKO-EMBASSY}记载，景初二年六月（238年），曹魏与倭女王国发生“倭女王卑弥呼遣使至魏，魏朝廷接受其朝贡并授予册封、铜镜等物”。这一节点可放在魏平定辽东后恢复经带方郡通向朝鲜半岛与海东的信息和朝贡渠道，倭地竞争者也需要外部承认的条件下理解；其后续影响被账本概括为魏的册封把倭地纳入其文书化外交视野，但不证明魏在列岛拥有直接统治；铜镜与考古发现的对应须逐件核对。就地方尺度而言，事件并不只属于朝廷或统帅：征发、道路、仓储、户籍、市场、寺院与家庭关系都可能决定命令如何抵达、战事如何延长，或接收何以在不同地区呈现不一样的速度。

本条列举的核心材料为《三国志·魏书·乌丸鲜卑东夷传》；《三国志·明帝纪》；《日本书纪》后出相关记载。这些材料较适合钉住事件的发生、官方表述或特定人物、地点，却不自动构成对全境人口、收入、兵数、信仰或动机的调查。账本的置信与材料提示是“遣使年月和魏朝廷记载可靠；邪马台国位置、使团路线、授物数量及其在倭地的政治含义有争议”；来源限度进一步说明“魏倭外交、带方郡与东亚海域交流研究”。因此，不能把条目中的政权名称、行政分类或战报修辞直接转译为固定族群和统一社会意志，也不能将制度宣布写成各地同步实施。

在叙事上，更稳妥的做法是把这一事件视为一条需要地方中介才能运作的链：中央的命令、地方官与精英的选择、运输与劳力的条件、以及普通家庭可能面对的迁徙、赋役或安全风险彼此相连。现有材料能够显示其中若干环节，却保留了大量沉默。承认这种沉默并非放弃解释，而是避免以一条编年记录覆盖不同区域和社群的差异。

账本条目\texttt{TK-0239-CAO-RUI-DEATH}记载，景初三年正月（239年），曹魏发生“曹叡去世，幼帝曹芳即位，曹爽与司马懿受遗诏辅政”。这一节点可放在曹叡无成年继承人，且晚年宫室营建与边境战事加重中枢压力，必须以辅政者维持军政连续的条件下理解；其后续影响被账本概括为曹芳朝形成曹爽与司马懿并立的权力结构，魏的皇帝继承问题转化为辅政集团争夺诏令、禁军和任官权。就地方尺度而言，事件并不只属于朝廷或统帅：征发、道路、仓储、户籍、市场、寺院与家庭关系都可能决定命令如何抵达、战事如何延长，或接收何以在不同地区呈现不一样的速度。

本条列举的核心材料为《三国志·明帝纪》；《三国志·齐王芳纪》；《三国志·曹爽传》；《资治通鉴》魏纪。这些材料较适合钉住事件的发生、官方表述或特定人物、地点，却不自动构成对全境人口、收入、兵数、信仰或动机的调查。账本的置信与材料提示是“卒年、继位和辅政安排可靠；遗诏文本、曹爽司马懿权力均衡及近侍作用有争议”；来源限度进一步说明“魏明帝身后辅政、幼主政治与曹爽司马懿研究”。因此，不能把条目中的政权名称、行政分类或战报修辞直接转译为固定族群和统一社会意志，也不能将制度宣布写成各地同步实施。

在叙事上，更稳妥的做法是把这一事件视为一条需要地方中介才能运作的链：中央的命令、地方官与精英的选择、运输与劳力的条件、以及普通家庭可能面对的迁徙、赋役或安全风险彼此相连。现有材料能够显示其中若干环节，却保留了大量沉默。承认这种沉默并非放弃解释，而是避免以一条编年记录覆盖不同区域和社群的差异。

账本条目\texttt{TK-0240-ZHENGSHI-STONE-CLASSICS}记载，正始年间（240年至249年），曹魏发生“魏廷在洛阳太学立三体石经，以古文、篆、隶并列校写经典”。这一节点可放在魏中枢需要借可公开校读的经文稳定官学与书写权威，回应汉末以来传本、字形和师法差异的条件下理解；其后续影响被账本概括为石经为魏的教育和合法性提供可见媒介；残石能校验部分文字，不能据此推定全国学校或经学解释已经统一。就地方尺度而言，事件并不只属于朝廷或统帅：征发、道路、仓储、户籍、市场、寺院与家庭关系都可能决定命令如何抵达、战事如何延长，或接收何以在不同地区呈现不一样的速度。

本条列举的核心材料为《三国志·齐王芳纪》；《晋书·卫恒传》；正始三体石经残石及洛阳太学遗址资料。这些材料较适合钉住事件的发生、官方表述或特定人物、地点，却不自动构成对全境人口、收入、兵数、信仰或动机的调查。账本的置信与材料提示是“立石经及其正始年代大体可靠；始刻年份、篇目、残石归属与教学使用范围有争议”；来源限度进一步说明“正始石经、经学定本与洛阳物质文化研究”。因此，不能把条目中的政权名称、行政分类或战报修辞直接转译为固定族群和统一社会意志，也不能将制度宣布写成各地同步实施。

在叙事上，更稳妥的做法是把这一事件视为一条需要地方中介才能运作的链：中央的命令、地方官与精英的选择、运输与劳力的条件、以及普通家庭可能面对的迁徙、赋役或安全风险彼此相连。现有材料能够显示其中若干环节，却保留了大量沉默。承认这种沉默并非放弃解释，而是避免以一条编年记录覆盖不同区域和社群的差异。

账本条目\texttt{TK-0241-SUN-DENG-DEATH}记载，赤乌四年（241年），孙吴发生“太子孙登去世，孙吴失去已被长期培养的成年继承人”。这一节点可放在孙登原可在孙权晚年分担政务并平衡江东官僚、宗室与将领，其死亡使继承安排重新开放的条件下理解；其后续影响被账本概括为孙权后来立孙和、封孙霸，储位竞争的条件形成；继承危机将削弱孙吴在淮南与荆州的决策稳定性。就地方尺度而言，事件并不只属于朝廷或统帅：征发、道路、仓储、户籍、市场、寺院与家庭关系都可能决定命令如何抵达、战事如何延长，或接收何以在不同地区呈现不一样的速度。

本条列举的核心材料为《三国志·孙登传》；《三国志·吴主传》；《建康实录》。这些材料较适合钉住事件的发生、官方表述或特定人物、地点，却不自动构成对全境人口、收入、兵数、信仰或动机的调查。账本的置信与材料提示是“卒年与太子身份可靠；孙登病因、其政治角色和各派反应有争议”；来源限度进一步说明“孙登之死、孙吴储位与宗室政治研究”。因此，不能把条目中的政权名称、行政分类或战报修辞直接转译为固定族群和统一社会意志，也不能将制度宣布写成各地同步实施。

在叙事上，更稳妥的做法是把这一事件视为一条需要地方中介才能运作的链：中央的命令、地方官与精英的选择、运输与劳力的条件、以及普通家庭可能面对的迁徙、赋役或安全风险彼此相连。现有材料能够显示其中若干环节，却保留了大量沉默。承认这种沉默并非放弃解释，而是避免以一条编年记录覆盖不同区域和社群的差异。

\subsection{信仰、知识与物质空间}

账本条目\texttt{TK-0242-SUN-HE-SUN-BA}记载，赤乌五年（242年），孙吴发生“孙和被立为太子，孙霸封鲁王，两个成年皇子各自形成宾客与官僚支持网络”。这一节点可放在孙登死亡后孙权需尽快确立继承人，却同时给予孙霸高位和可观的政治资源的条件下理解；其后续影响被账本概括为太子与鲁王并立使中枢官员被迫选边，孙权晚年的人事与军事决策将受到持续的储位冲突牵制。就地方尺度而言，事件并不只属于朝廷或统帅：征发、道路、仓储、户籍、市场、寺院与家庭关系都可能决定命令如何抵达、战事如何延长，或接收何以在不同地区呈现不一样的速度。

本条列举的核心材料为《三国志·孙和传》；《三国志·孙霸传》裴松之注引《吴录》；《三国志·吴主传》。这些材料较适合钉住事件的发生、官方表述或特定人物、地点，却不自动构成对全境人口、收入、兵数、信仰或动机的调查。账本的置信与材料提示是“册立和封王可靠；两府权力、党附规模和冲突起点有争议”；来源限度进一步说明“孙吴二宫之争、宗室宾客与晚年政治研究”。因此，不能把条目中的政权名称、行政分类或战报修辞直接转译为固定族群和统一社会意志，也不能将制度宣布写成各地同步实施。

在叙事上，更稳妥的做法是把这一事件视为一条需要地方中介才能运作的链：中央的命令、地方官与精英的选择、运输与劳力的条件、以及普通家庭可能面对的迁徙、赋役或安全风险彼此相连。现有材料能够显示其中若干环节，却保留了大量沉默。承认这种沉默并非放弃解释，而是避免以一条编年记录覆盖不同区域和社群的差异。

账本条目\texttt{TK-0244-CAO-SHUANG-XINGSHI}记载，正始五年（244年），曹魏与蜀汉发生“曹爽率军攻蜀至兴势，遭王平等据险防守而撤退”。这一节点可放在曹爽希望以对蜀军功巩固辅政威望，魏军却须沿狭窄山道穿越汉中外围的险隘的条件下理解；其后续影响被账本概括为魏军未能取得汉中，曹爽的军事声望受损；蜀汉以较少兵力守险的经验影响其后续防御部署。就地方尺度而言，事件并不只属于朝廷或统帅：征发、道路、仓储、户籍、市场、寺院与家庭关系都可能决定命令如何抵达、战事如何延长，或接收何以在不同地区呈现不一样的速度。

本条列举的核心材料为《三国志·曹爽传》；《三国志·王平传》；《三国志·费祎传》；《资治通鉴》魏纪。这些材料较适合钉住事件的发生、官方表述或特定人物、地点，却不自动构成对全境人口、收入、兵数、信仰或动机的调查。账本的置信与材料提示是“出师、受阻与撤军可靠；行军路线、伤亡、曹爽亲自指挥程度和费祎部署有争议”；来源限度进一步说明“兴势之役、曹爽伐蜀与汉中山地防御研究”。因此，不能把条目中的政权名称、行政分类或战报修辞直接转译为固定族群和统一社会意志，也不能将制度宣布写成各地同步实施。

在叙事上，更稳妥的做法是把这一事件视为一条需要地方中介才能运作的链：中央的命令、地方官与精英的选择、运输与劳力的条件、以及普通家庭可能面对的迁徙、赋役或安全风险彼此相连。现有材料能够显示其中若干环节，却保留了大量沉默。承认这种沉默并非放弃解释，而是避免以一条编年记录覆盖不同区域和社群的差异。

账本条目\texttt{TK-0244-GOGURYEO-CAMPAIGN}记载，正始五年至六年（244年至245年），曹魏与高句丽发生“毌丘俭两次进攻高句丽，攻破丸都城，高句丽王东川王避走”。这一节点可放在魏完成辽东征服后试图压制高句丽对玄菟及东部交通的影响，高句丽也因边境扩张与魏军冲突的条件下理解；其后续影响被账本概括为魏军短期深入山地并破都城，但未建立长期直接统治；中韩文献和丸都山城考古须并列校读。就地方尺度而言，事件并不只属于朝廷或统帅：征发、道路、仓储、户籍、市场、寺院与家庭关系都可能决定命令如何抵达、战事如何延长，或接收何以在不同地区呈现不一样的速度。

本条列举的核心材料为《三国志·毌丘俭传》；《三国志·魏书·乌丸鲜卑东夷传》；《三国史记·高句丽本纪》。这些材料较适合钉住事件的发生、官方表述或特定人物、地点，却不自动构成对全境人口、收入、兵数、信仰或动机的调查。账本的置信与材料提示是“魏军出征与破城大体可靠；行军路线、城市位置、战果和双方纪年的换算有争议”；来源限度进一步说明“魏与高句丽战争、丸都城考古与东北边疆研究”。因此，不能把条目中的政权名称、行政分类或战报修辞直接转译为固定族群和统一社会意志，也不能将制度宣布写成各地同步实施。

在叙事上，更稳妥的做法是把这一事件视为一条需要地方中介才能运作的链：中央的命令、地方官与精英的选择、运输与劳力的条件、以及普通家庭可能面对的迁徙、赋役或安全风险彼此相连。现有材料能够显示其中若干环节，却保留了大量沉默。承认这种沉默并非放弃解释，而是避免以一条编年记录覆盖不同区域和社群的差异。

账本条目\texttt{TK-0245-LU-XUN-DEATH}记载，赤乌八年（245年），孙吴发生“丞相陆逊去世，其与孙权围绕太子孙和、鲁王孙霸的谏争成为后世叙述焦点”。这一节点可放在二宫党争侵入军政人事，陆逊以丞相和武昌方面重臣身份干预储位，触及孙权晚年的决断的条件下理解；其后续影响被账本概括为孙吴失去兼具荆州军政经验的重臣，储位冲突更难调停；后出书信需要辨析其传抄和政治立场。就地方尺度而言，事件并不只属于朝廷或统帅：征发、道路、仓储、户籍、市场、寺院与家庭关系都可能决定命令如何抵达、战事如何延长，或接收何以在不同地区呈现不一样的速度。

本条列举的核心材料为《三国志·陆逊传》；《三国志·孙和传》；《三国志·陆逊传》裴松之注引《吴书》。这些材料较适合钉住事件的发生、官方表述或特定人物、地点，却不自动构成对全境人口、收入、兵数、信仰或动机的调查。账本的置信与材料提示是“卒年可靠；陆逊受责、病死与储位争论之间的因果及书信真实性有争议”；来源限度进一步说明“陆逊之死、二宫之争与孙吴重臣政治研究”。因此，不能把条目中的政权名称、行政分类或战报修辞直接转译为固定族群和统一社会意志，也不能将制度宣布写成各地同步实施。

在叙事上，更稳妥的做法是把这一事件视为一条需要地方中介才能运作的链：中央的命令、地方官与精英的选择、运输与劳力的条件、以及普通家庭可能面对的迁徙、赋役或安全风险彼此相连。现有材料能够显示其中若干环节，却保留了大量沉默。承认这种沉默并非放弃解释，而是避免以一条编年记录覆盖不同区域和社群的差异。

账本条目\texttt{TK-0247-JIANG-WEI-LONGXI}记载，延熙十年（247年），蜀汉、曹魏与陇西羌胡集团发生“姜维趁陇西羌胡起事出兵，魏将郭淮等应战，蜀汉重新将北伐重心推向陇右”。这一节点可放在蜀汉在费祎主持下保持边境防御，陇西部众与魏郡县的冲突提供姜维介入和牵制魏军的机会的条件下理解；其后续影响被账本概括为姜维成为蜀汉北方军事的关键行动者；将羌胡行动直接归入蜀汉盟军会掩盖其自身的地方政治目标。就地方尺度而言，事件并不只属于朝廷或统帅：征发、道路、仓储、户籍、市场、寺院与家庭关系都可能决定命令如何抵达、战事如何延长，或接收何以在不同地区呈现不一样的速度。

本条列举的核心材料为《三国志·姜维传》；《三国志·郭淮传》；《资治通鉴》魏纪。这些材料较适合钉住事件的发生、官方表述或特定人物、地点，却不自动构成对全境人口、收入、兵数、信仰或动机的调查。账本的置信与材料提示是“出兵和魏蜀交战可靠；羌胡各部关系、地方响应和姜维战略自主性有争议”；来源限度进一步说明“姜维早期北伐、陇西羌胡与蜀魏边疆研究”。因此，不能把条目中的政权名称、行政分类或战报修辞直接转译为固定族群和统一社会意志，也不能将制度宣布写成各地同步实施。

在叙事上，更稳妥的做法是把这一事件视为一条需要地方中介才能运作的链：中央的命令、地方官与精英的选择、运输与劳力的条件、以及普通家庭可能面对的迁徙、赋役或安全风险彼此相连。现有材料能够显示其中若干环节，却保留了大量沉默。承认这种沉默并非放弃解释，而是避免以一条编年记录覆盖不同区域和社群的差异。

账本条目\texttt{TK-0249-GAOPING-TOMBS}记载，正始十年正月（249年），曹魏发生“司马懿发动高平陵政变，控制洛阳并罢黜曹爽集团”。这一节点可放在曹爽长期排挤司马懿并掌握皇帝近侍与军政任命，曹芳出谒高平陵时暴露洛阳城内权力空隙的条件下理解；其后续影响被账本概括为司马氏控制魏的中枢和人事，曹魏皇帝继续在位但辅政体制转为司马氏主导，为后来的禅代铺设制度通道。就地方尺度而言，事件并不只属于朝廷或统帅：征发、道路、仓储、户籍、市场、寺院与家庭关系都可能决定命令如何抵达、战事如何延长，或接收何以在不同地区呈现不一样的速度。

本条列举的核心材料为《三国志·曹爽传》；《三国志·司马懿传》；《三国志·齐王芳纪》；《资治通鉴》魏纪。这些材料较适合钉住事件的发生、官方表述或特定人物、地点，却不自动构成对全境人口、收入、兵数、信仰或动机的调查。账本的置信与材料提示是“政变、罢黜和曹爽随后被杀可靠；政变时的军队部署、诏令真伪及司马懿承诺是否构成欺罔有争议”；来源限度进一步说明“高平陵政变、魏都禁军与辅政权力研究”。因此，不能把条目中的政权名称、行政分类或战报修辞直接转译为固定族群和统一社会意志，也不能将制度宣布写成各地同步实施。

在叙事上，更稳妥的做法是把这一事件视为一条需要地方中介才能运作的链：中央的命令、地方官与精英的选择、运输与劳力的条件、以及普通家庭可能面对的迁徙、赋役或安全风险彼此相连。现有材料能够显示其中若干环节，却保留了大量沉默。承认这种沉默并非放弃解释，而是避免以一条编年记录覆盖不同区域和社群的差异。

账本条目\texttt{TK-0250-SUN-WU-SUCCESSION-CRISIS}记载，赤乌十三年（250年），孙吴发生“孙权废太子孙和、赐死鲁王孙霸，改立幼子孙亮为太子，二宫之争以强制清洗收束”。这一节点可放在孙和与孙霸各自形成的宾客集团长期侵蚀朝廷，孙权既不能容忍持续分裂，也缺少能兼得各方认可的成年继承人的条件下理解；其后续影响被账本概括为大量官员被诛逐，孙亮继承使孙吴再次面对幼主与权臣问题；事件不能被化约为单纯的家族伦理冲突。就地方尺度而言，事件并不只属于朝廷或统帅：征发、道路、仓储、户籍、市场、寺院与家庭关系都可能决定命令如何抵达、战事如何延长，或接收何以在不同地区呈现不一样的速度。

本条列举的核心材料为《三国志·孙和传》；《三国志·孙亮传》；《三国志·吴主传》；《资治通鉴》魏纪。这些材料较适合钉住事件的发生、官方表述或特定人物、地点，却不自动构成对全境人口、收入、兵数、信仰或动机的调查。账本的置信与材料提示是“废立和孙霸死亡可靠；参与者罪责、孙和处置细节和孙亮被立的决策过程有争议”；来源限度进一步说明“孙吴二宫之争结局、储位清洗与幼主继承研究”。因此，不能把条目中的政权名称、行政分类或战报修辞直接转译为固定族群和统一社会意志，也不能将制度宣布写成各地同步实施。

在叙事上，更稳妥的做法是把这一事件视为一条需要地方中介才能运作的链：中央的命令、地方官与精英的选择、运输与劳力的条件、以及普通家庭可能面对的迁徙、赋役或安全风险彼此相连。现有材料能够显示其中若干环节，却保留了大量沉默。承认这种沉默并非放弃解释，而是避免以一条编年记录覆盖不同区域和社群的差异。

账本条目\texttt{TK-0251-WANG-LING-REBELLION}记载，嘉平三年（251年），曹魏发生“太尉王凌谋立楚王曹彪以反司马懿，事泄后投降自杀，曹彪被迫自尽”。这一节点可放在高平陵后司马氏迅速控制朝廷，部分魏臣和宗室将淮南军政资源视为恢复曹氏实权的可能依托的条件下理解；其后续影响被账本概括为司马懿清除首个大型反司马氏网络，楚王一支被削弱；密谋书信和告发材料须注意胜方档案的选择性。就地方尺度而言，事件并不只属于朝廷或统帅：征发、道路、仓储、户籍、市场、寺院与家庭关系都可能决定命令如何抵达、战事如何延长，或接收何以在不同地区呈现不一样的速度。

本条列举的核心材料为《三国志·王凌传》；《三国志·齐王芳纪》；《三国志·王凌传》裴松之注引《魏略》。这些材料较适合钉住事件的发生、官方表述或特定人物、地点，却不自动构成对全境人口、收入、兵数、信仰或动机的调查。账本的置信与材料提示是“谋反、处置和死亡可靠；密谋范围、诏书真伪与淮南地方支持程度有争议”；来源限度进一步说明“王凌之变、淮南与曹魏宗室反司马氏研究”。因此，不能把条目中的政权名称、行政分类或战报修辞直接转译为固定族群和统一社会意志，也不能将制度宣布写成各地同步实施。

在叙事上，更稳妥的做法是把这一事件视为一条需要地方中介才能运作的链：中央的命令、地方官与精英的选择、运输与劳力的条件、以及普通家庭可能面对的迁徙、赋役或安全风险彼此相连。现有材料能够显示其中若干环节，却保留了大量沉默。承认这种沉默并非放弃解释，而是避免以一条编年记录覆盖不同区域和社群的差异。

账本条目\texttt{TK-0251-SIMA-YI-DEATH}记载，嘉平三年八月（251年），曹魏发生“司马懿去世，司马师继承其父的中军与辅政地位”。这一节点可放在司马懿已借高平陵与王凌案重塑中枢，司马氏需要让军政资源在家族内连续转移而不触发新的反弹的条件下理解；其后续影响被账本概括为司马师成为魏朝实际主导者，司马氏的家族化辅政进一步削弱皇帝对禁军和任官的控制。就地方尺度而言，事件并不只属于朝廷或统帅：征发、道路、仓储、户籍、市场、寺院与家庭关系都可能决定命令如何抵达、战事如何延长，或接收何以在不同地区呈现不一样的速度。

本条列举的核心材料为《三国志·司马懿传》；《三国志·司马师传》；《资治通鉴》魏纪。这些材料较适合钉住事件的发生、官方表述或特定人物、地点，却不自动构成对全境人口、收入、兵数、信仰或动机的调查。账本的置信与材料提示是“卒年和司马师继掌权力可靠；继承的具体官职、军权交接和曹芳的自主空间有争议”；来源限度进一步说明“司马懿身后继承、司马师与魏国权力结构研究”。因此，不能把条目中的政权名称、行政分类或战报修辞直接转译为固定族群和统一社会意志，也不能将制度宣布写成各地同步实施。

在叙事上，更稳妥的做法是把这一事件视为一条需要地方中介才能运作的链：中央的命令、地方官与精英的选择、运输与劳力的条件、以及普通家庭可能面对的迁徙、赋役或安全风险彼此相连。现有材料能够显示其中若干环节，却保留了大量沉默。承认这种沉默并非放弃解释，而是避免以一条编年记录覆盖不同区域和社群的差异。

账本条目\texttt{TK-0252-SUN-QUAN-DEATH}记载，太元二年（252年），孙吴发生“孙权去世，太子孙亮继位，诸葛恪受遗诏辅政”。这一节点可放在二宫之争已清洗大批重臣，孙亮年幼，孙权只能依靠外戚性重臣和宗室军人维持中枢与边防的条件下理解；其后续影响被账本概括为孙吴进入诸葛恪主政期，辅政集团的军事胜负将直接影响幼帝和江东官僚对其合法性的判断。就地方尺度而言，事件并不只属于朝廷或统帅：征发、道路、仓储、户籍、市场、寺院与家庭关系都可能决定命令如何抵达、战事如何延长，或接收何以在不同地区呈现不一样的速度。

本条列举的核心材料为《三国志·吴主传》；《三国志·孙亮传》；《三国志·诸葛恪传》。这些材料较适合钉住事件的发生、官方表述或特定人物、地点，却不自动构成对全境人口、收入、兵数、信仰或动机的调查。账本的置信与材料提示是“卒年、继位和辅政事实可靠；遗诏内容、孙峻等宗室军人的位置和诸葛恪实际权限有争议”；来源限度进一步说明“孙权身后继承、孙亮幼主与诸葛恪辅政研究”。因此，不能把条目中的政权名称、行政分类或战报修辞直接转译为固定族群和统一社会意志，也不能将制度宣布写成各地同步实施。

在叙事上，更稳妥的做法是把这一事件视为一条需要地方中介才能运作的链：中央的命令、地方官与精英的选择、运输与劳力的条件、以及普通家庭可能面对的迁徙、赋役或安全风险彼此相连。现有材料能够显示其中若干环节，却保留了大量沉默。承认这种沉默并非放弃解释，而是避免以一条编年记录覆盖不同区域和社群的差异。

账本条目\texttt{TK-0252-DONGXING}记载，建兴元年（252年），孙吴与曹魏发生“诸葛恪在东兴堤击退魏军，吴军利用水陆地形保住濡须方向”。这一节点可放在孙权死后魏试图以大军试探幼主政权，东兴堤和湖泊水网成为双方争夺的前沿据点的条件下理解；其后续影响被账本概括为孙吴取得诸葛恪辅政后的重要胜利，短期提振江东士气；传记的突击叙事不宜直接换算为精确战损。就地方尺度而言，事件并不只属于朝廷或统帅：征发、道路、仓储、户籍、市场、寺院与家庭关系都可能决定命令如何抵达、战事如何延长，或接收何以在不同地区呈现不一样的速度。

本条列举的核心材料为《三国志·诸葛恪传》；《三国志·丁奉传》；《三国志·齐王芳纪》；《资治通鉴》魏纪。这些材料较适合钉住事件的发生、官方表述或特定人物、地点，却不自动构成对全境人口、收入、兵数、信仰或动机的调查。账本的置信与材料提示是“战役胜负可靠；堤坝位置、参战兵力、丁奉突击细节和伤亡数字有争议”；来源限度进一步说明“东兴之战、水利地形与魏吴淮南攻防研究”。因此，不能把条目中的政权名称、行政分类或战报修辞直接转译为固定族群和统一社会意志，也不能将制度宣布写成各地同步实施。

在叙事上，更稳妥的做法是把这一事件视为一条需要地方中介才能运作的链：中央的命令、地方官与精英的选择、运输与劳力的条件、以及普通家庭可能面对的迁徙、赋役或安全风险彼此相连。现有材料能够显示其中若干环节，却保留了大量沉默。承认这种沉默并非放弃解释，而是避免以一条编年记录覆盖不同区域和社群的差异。
